\documentclass[../../../include/open-logic-section]{subfiles}

\begin{document}

\olfileid{sth}{card-arithmetic}{fix}
\olsection{Titik Ajeg $\aleph$}

Ing \olref[spine][]{chap}, kita ngandhakake yèn Aksioma Panggantian
mbutuhaké pambenaran, amarga aksioma iki ndadèkaké hirarki
dadi dhuwur banget. Sawisé nyinaoni aritmetika kardinal, kita bisa
mènèhi gambaran cilik bab dhuwuré hirarki kasebut.

Cetha yèn $0 < \aleph_0$, $1 < \aleph_1$, lan $2 < \aleph_2$\ldots;
malah béda ukurané saya \emph{gedhé} ing saben langkah. Mula ana
godhaan kanggo ngira yèn $\kappa< \aleph_\kappa$ kanggo saben
ordinal $\kappa$.

Nanging pangira iki \emph{klèru}, miturut $\ZFC$. Malah kita bisa
mbuktekaké anané \emph{titik ajeg $\aleph$}, yaiku kardinal
$\kappa$ kang nyukupi $\kappa=\aleph_\kappa$.

\begin{prop}\ollabel{alephfixed}
Ana titik ajeg $\aleph$.
\end{prop}

\begin{proof}
Kanthi rekursi, tetepna:
\begin{align*}
	\kappa_0 &= 0\\
	\kappa_{n+1} &= \aleph_{\kappa_n}\\
	\kappa&= \bigcup_{n < \omega}\kappa_n
\end{align*}
Saiki $\kappa$ iku kardinal miturut
\olref[cardinals][classing]{unioncardinalscardinal}. Sabanjuré:
\[
	\kappa= \bigcup_{n < \omega} \kappa_{n+1} =
	\bigcup_{n < \omega}\aleph_{\kappa_n} =
	\bigcup_{\alpha < \kappa}\aleph_\alpha = \aleph_\kappa
\]
%	Miturut konstruksi, $\kappa$ iku kardinal paling cilik kang luwih gedhé
%	imbang saben $\kappa_n$. Mula $\aleph_\kappa$ iku kardinal paling cilik
%	kang luwih gedhé tinimbang saben $\aleph_{\kappa_n}$, lan mula uga
%	luwih gedhé tinimbang saben $\kappa_{n+1}$. Nanging $\kappa$ uga
%	kardinal paling cilik kang luwih gedhé tinimbang saben
%	$\kappa_{n+1} = \aleph_{\kappa_n}$. Dadi $\kappa=
%	\aleph_\kappa$.
\end{proof}

Boolos naté nulis artikel bab titik ajeg $\aleph$ kang mentas kita
konstruksi. Sawisé nyebut anané $\kappa$ ing wiwitan artikelé,
dhèwèké ngandika:
\begin{quote}
	[$\kappa$ iku] cacah kang \emph{gedhé banget}, miturut pangertèné
	wong kang durung naté sinau teori himpunan; nganti, miturutku,
	kahanan iku ndadèkaké bebeneré sembarang téori kang negesaké
	anaké paling ora $\kappa$ obyèk dadi pitakonan.
	\citep[kaca~257]{Boolos2000}
\end{quote}
Ing pungkasan artikelé dhèwèké takon:
\begin{quote}
	[apa] kita curiga yèn, apa waé kahanané ing wiwitan crita,
	manawa wis tekan kéné rodhané mung muter tanpa tujuan lan
	kita ora lagi ngrungokaké gambaran bab samubarang kang nyata?
	\citep[kaca~268]{Boolos2000}
\end{quote}
Yèn pancèn kita wis ngluwihi ``samubarang kang nyata'', mula
panyalahé kudu kita tujokaké langsung marang Aksioma Panggantian.
Aksioma iki kang ndadèkaké bukti mau bisa lumaku. Yèn mangkono,
Boolos mesthiné kudu nliti maneh pratelané sawatara dasawarsa
sadurungé yèn Panggantian ora nduwèni akibat kang ``ora dikarepaké''
(delengen \olref[replacement][extrinsic]{sec}).

Nanging apa anané $\kappa$ pancèn ala? Ing kéné bisa migunani
nyawang \emph{paradoks Tristram Shandy} saka Russell. Tristram Shandy
nyathet uripé ing buku harian, nanging dhèwèké mbutuhaké setaun kanggo
nyathet sedina waé. Saben taun liwati, cathetané saya keri: sawisé
setaun dhèwèké lagi nyathet sedina lan isih ana 364 dina kang durung
kacathet; sawisé rong taun, lagi rong dina kang kacathet lan 728 dina
durung kacathet; sawisé telung taun, lagi telung dina kang kacathet
lan 1092 dina durung kacathet \dots\footnote{Taun kabisat ora diétung.}
Nanging yèn Tristram \emph{ora bisa mati}, dhèwèké bakal kasil
nyathet saben dina, amarga dina kaping $n$ bakal kacathet ing taun
kaping $n$ uripé. Mula ``ing pungkasaning wektu'' dhèwèké bakal
nduwèni buku harian kang jangkep.

Saiki, apa bédané karo pamikiran yèn $\alpha$ luwih cilik tinimbang
$\aleph_\alpha$---malah saya cilik banget---nganti tekan $\kappa$,
nalika rong sisih mau padha tekan titik kang padha lan $\kappa
= \aleph_\kappa$?

Kita sisihaké dhisik pitakonan iku. Yèn kita nampa $\ZFC$, ayo
dipungkasi kanthi konstruksi titik ajeg kang luwih nyenengaké.
Telung asil sabanjuré nerangaké, miturut pangertèn intuitif, yèn
ana titik kang ora sepele nalika amba lan dhuwuré hirarki padha:

\begin{prop}\ollabel{bethfixed}
Ana titik ajeg $\beth$, yaiku $\kappa$ kang nyukupi $\kappa=
\beth_\kappa$.
\end{prop}

\begin{proof}
Kaya ing \olref{alephfixed}, nanging ``$\beth$'' ngganti ``$\aleph$''.
\end{proof}

\begin{prop}\ollabel{stagesize}
$\card{V_{\omega+\alpha}} = \beth_{\alpha}$. Yèn $\omega \ordtimes
\omega \leq \alpha$, mula $\card{V_\alpha} = \beth_\alpha$.
\end{prop}

\begin{proof}
Pratelan kapisan kabukten kanthi induksi transfinit kang prasaja.
Pratelan kapindho katut amarga yèn $\omega \ordtimes \omega \leq \alpha$,
mula $\omega + \alpha = \alpha$. Kanggo mbuktekaké iki, kita migunakaké
kasunyatan aritmetika ordinal saka \olref[ord-arithmetic][]{chap}.
Wiwitané, gatekna $\omega \ordtimes \omega = \omega \ordtimes (1 \ordplus \omega) =
(\omega  \ordtimes 1) \ordplus (\omega\ordtimes\omega) = \omega
\ordplus (\omega \ordtimes \omega)$. Saiki yèn $\omega \ordtimes \omega
\leq \alpha$, yaiku $\alpha = (\omega\ordtimes\omega) \ordplus \beta$
kanggo sawijining $\beta$, mula $\omega \ordplus \alpha = \omega \ordplus
((\omega \ordtimes \omega) \ordplus \beta) = (\omega \ordplus (\omega
\ordtimes \omega)) \ordplus \beta = (\omega \ordtimes \omega) \ordplus
\beta = \alpha$.
\end{proof}

\begin{cor}
Ana $\kappa$ kang nyukupi $\card{V_\kappa} = \kappa$.
\end{cor}

\begin{proof}
Jupuken $\kappa$ titik ajeg $\beth$ kaya kang diwènèhaké déning
\olref{bethfixed}. Cetha $\omega \ordtimes \omega < \kappa$.
Mula $\card{V_\kappa} = \beth_\kappa= \kappa$ miturut
\olref{stagesize}.
\end{proof}

Cacahé tahap ing sangisoré $V_\kappa$ padha karo cacahé !!{element}s
ing $V_\kappa$. Mula miturut pangertèn intuitif, $V_\kappa$ padha
amba lan dhuwuré. Iki ngélingaké marang Tristram Shandy: kita pindhah
saka siji tahap menyang tahap sabanjuré kanthi njupuk
\emph{himpunan pangkat}, saéngga hirarki dadi \emph{luwih}
gedhé ing saben langkah. Nanging ``ing pungkasan'', yaiku ing tahap
$\kappa$, amba hirarki mau nyusul dhuwuré.

Wong bisa takon: \emph{Sepira kerepé amba hirarki padha karo
dhuwuré?} Wangsulané: \emph{Padha kerepé karo anané ordinal.}
Nanging iki perlu dijlentrehaké sethithik.

Kita dhéfinisikaké suku $\tau$ kaya mangkéné. Kanggo sembarang $A$,
tetepna:
\begin{align*}
	\tau_0(A) & \defis \cardsucc{\card{A}}\\
	\tau_{n+1}(A) & \defis \beth_{\tau_n(A)}\\
	\tau(A) & \defis \bigcup_{n < \omega}\tau_n(A)
\intertext{Kaya ing \olref{bethfixed}, $\tau(A)$ iku titik ajeg
$\beth$ kanggo saben $A$, lan mesthi $\card{A} < \tau(A)$.
Saiki gatekna dhéfinisi rekursif iki:}
	W_0 &\defis \tau(0)\\
	W_{\alpha + 1} & \defis \tau(W_\alpha)\\
	W_\alpha & \defis \bigcup_{\beta < \alpha} W_\beta
	\text{, nalika $\alpha$ ordinal limit}
\end{align*}
Konstruksi iki didhéfinisikaké kanggo saben ordinal. Ing tahap limit,
gabungan runtutan kang mundhak saka titik ajeg bet tetep titik ajeg bet.
Mula $W$ iku ``!!a{injection}'' kelas saka ordinal menyang titik ajeg
$\beth$. Lan kaya sadurungé, $V_{W_\alpha}$ padha amba lan dhuwuré
kanggo saben $\alpha$.

\end{document}
